\section{mamabench: A Scope-Filtered QA Benchmark}
\label{sec:mamabench}

mamabench evaluates how well a model answers the questions a nurse-midwife
asks---those spanning maternal, neonatal, child, and reproductive health---in
formats from multiple-choice to free-text clinical scenarios. We build it by
assembling such questions from seven existing, expert-authored sources and
filtering each to that scope. This poses three construction problems: assembling
and normalising the sources into one scorable schema, defining and applying the
scope reliably, and giving consumers a way to trust the LLM judge that scores the
open-ended items. We take each in turn.

\subsection{Source assembly}

The first decision is to assemble, not author: writing and validating
maternal-health questions from scratch needs clinicians, whereas the seven
sources already carry expert curation---medical-licensing boards, a pan-African
physician panel, Kenyan clinicians, and OpenAI's physician annotators---so the
benchmark's contribution is scoping and uniform packaging, not de-novo clinical
authoring.
Each source is normalised by a dedicated adapter into one canonical schema and
sorted into one of three tracks (Table~\ref{tab:mamabench-sources}):

\begin{itemize}
  \item \textbf{Multiple-choice} (23{,}241 items), for breadth of clinical
  knowledge: the OBGYN and paediatrics subset of MedMCQA~\cite{pal2022medmcqa},
  the in-scope portion of MedQA-USMLE~\cite{jin2021medqa}, and the
  multiple-choice questions of AfriMed-QA~\cite{olatunji2025afrimedqa}. This
  track serves breadth; the open-ended tracks below carry the weight.
  \item \textbf{Short-answer} (369 items), closest to the deployment setting:
  312 nurse-written primary-care cases~\cite{kenya2025vignettes}, the
  short-answer questions of AfriMed-QA~\cite{olatunji2025afrimedqa}, and a small
  set of expert-crafted items designed to expose model
  errors~\cite{whb2025}.
  \item \textbf{Rubric-graded} (2{,}339 items): the in-scope portion of
  HealthBench~\cite{arora2025healthbench}, whose items are scored against
  physician-written rubric criteria rather than a single gold answer.
\end{itemize}

\begin{table}[t]
\centering
\small
\setlength{\tabcolsep}{6pt}
\renewcommand{\arraystretch}{1.12}
\caption{The seven sources of mamabench, by track. \emph{Yield} is the share of
a source's upstream pool that our scope classifier (\S\ref{sec:mamabench-scope})
keeps as in-scope; a dash marks sources that arrive already scoped to the
domain---MedMCQA and AfriMed-QA by their subject or specialty labels, the
Women's Health Benchmark by expert authoring---and so bypass the classifier
entirely.}
\label{tab:mamabench-sources}
\begin{tabular}{@{}llrr@{}}
\toprule
Track & Source & Items & Yield \\
\midrule
\multirow{3}{*}{Multiple-choice}
 & MedMCQA (OBGYN \& paediatrics) & 18{,}508 & -- \\
 & MedQA-USMLE                    &  4{,}199 & 29.2\% \\
 & AfriMed-QA                     &     534  & -- \\
\midrule
\multirow{3}{*}{Short-answer}
 & Kenya Clinical Vignettes       &     312  & 61.5\% \\
 & AfriMed-QA (short-answer)      &      37  & -- \\
 & Women's Health Benchmark       &      20  & -- \\
\midrule
\multirow{3}{*}{Rubric-graded}
 & HealthBench (oss\_eval)        &  1{,}209 & 24.2\% \\
 & HealthBench (consensus)        &     872  & 23.8\% \\
 & HealthBench (hard)             &     258  & 25.9\% \\
\midrule
\multicolumn{2}{@{}l}{\textbf{Total}} & \textbf{25{,}949} & \\
\bottomrule
\end{tabular}
\end{table}

Normalisation is mostly mechanical, but one source needs more than that: the
AfriMed-QA multiple-choice items carry quality issues the others do not, and
there we draw the line at \emph{scorability}. Multiple-choice items are scored by
the answer's letter, and we drop only what that scoring cannot represent: items
with more than one correct option (a single answer index cannot encode them),
and items whose correct-answer text is duplicated under another letter---if the
same option text sits at two positions, a model could pick the right text yet be
marked wrong for reporting the duplicate letter. Cosmetic issues that do not
affect scoring (embedded
option-letter prefixes, non-ASCII typography, terse stems) are left untouched
and documented in the manifest. The principle is that the benchmark may refuse
an item it cannot score, but must never change what an item means.

\subsection{Defining scope with an LLM classifier}
\label{sec:mamabench-scope}

The central construction decision is how to decide whether an item is in scope.
We define scope as four positive categories---\textsc{maternal},
\textsc{neonatal}, \textsc{child health}, and
\textsc{sexual and reproductive health} (SRH)---plus \textsc{none}, and we
classify
each item by its \emph{primary clinical concept, not the patient's
demographics}. A question about managing a wrist fracture in a patient
who happens to be pregnant is \textsc{none}---the fracture is the concept, the
pregnancy incidental---whereas postpartum depression is \textsc{maternal},
because here the postpartum context \emph{is} the concept. The classifier is
\texttt{Qwen3.6-27B-FP8},\footnote{Model card:
\url{https://huggingface.co/Qwen/Qwen3.6-27B-FP8}.} run
with reasoning enabled at temperature 0; it applies this rule to every candidate
item and emits a structured verdict with its reasoning retained for audit.

Trusting an LLM to draw the scope boundary needs evidence, and we offer two
checks. First, \textbf{cross-model agreement}: re-classifying the HealthBench
\texttt{oss\_eval} items with the much larger
\texttt{Qwen3.5-397B-A17B-FP8}\footnote{Model card:
\url{https://huggingface.co/Qwen/Qwen3.5-397B-A17B-FP8}.}---a stronger model from
the same family, so this checks self-consistency more than
independence---agrees with the 27B on
98.1\% of the 4{,}988 jointly classified items, and per-category counts differ
by under one percent; the handful of items on which the 27B fails to converge
are all \textsc{none} under the 397B, so they would be filtered out either way
and the net effect on in-scope counts is zero. Second, \textbf{parity against
prior labels}: the Kenya vignettes arrived with their own classifier labels,
and ours agree with them on 86.8\%, with the disagreements dominated by our rule
being deliberately stricter---rejecting items where only the patient's being
female, rather than the clinical concept, suggested OBGYN scope.

The per-source yield in Table~\ref{tab:mamabench-sources} is itself a useful
signal. The Kenyan primary-care vignettes are 61.5\% in scope, against
24--29\% for the general benchmarks: the maternal-health sources
concentrate exactly where the deployment lives, while the general sets
contribute a smaller, scope-filtered slice.

\subsection{A judge-calibration set, re-scoped from HealthBench}
\label{sec:mamabench-calibration}

The rubric-graded track is scored by an LLM judge, which raises the obvious
question of whether the judge can be trusted. Rather than answer that once for
our own judge, we ship the means to answer it for \emph{any} judge: a
calibration side-file of 6{,}853 (prompt, response, criterion) triples carrying
physicians' met / not-met labels. We want to be precise about provenance, since
it is easy to overclaim here. \emph{We did not produce these labels.} They are
the physician annotations HealthBench released for vetting rubric graders---its
meta-evaluation set~\cite{arora2025healthbench}; our only role is to re-scope
that set to this domain---keeping the triples whose prompts the scope classifier
(\S\ref{sec:mamabench-scope}) places in the four positive categories
(\textsc{child health} 3{,}239, \textsc{maternal} 2{,}284, \textsc{srh} 942,
\textsc{neonatal} 388, over 872 prompts)---and to package them as a ready-to-use
calibration file alongside the benchmark. A consumer runs a candidate judge over
these triples, compares its labels to the physicians', and decides whether to
trust it before scoring the benchmark; the companion system paper selects its
rubric judge in exactly this way~\cite{ren2026mamai}. One caveat: these 872
prompts \emph{are} the HealthBench-consensus rubric track's prompts---HealthBench
collected the side-file's physician labels on that same set---so calibrating a
judge here and then scoring the consensus track reuses the same prompts;
calibrate on the side-file, but read consensus-track scores with that overlap in
mind.
